\label{sec:introduction}

At a 2021 public hearing of North Carolina's redistricting commission, a resident
of Raleigh's Cameron Park neighborhood testified that splitting her historic
district across two congressional seats would deprive the neighborhood's
long-tenured homeowners of a coherent representative voice on federal preservation
funding, housing finance policy, and historic tax credit renewal.
Her testimony was noted but could not be quantified.
The commission had no tool that could measure whether Cameron Park and the
adjacent Five Points neighborhood shared a community bond that the Oakwood
tract across the freeway did not.

This paper provides that measurement.

The M-track community character framework, introduced in M.0~\citep{m0framework},
operationalizes diverse community character signals as edge weights in the
census-tract adjacency graph used by \textsc{bisect}.
Tracts whose character vectors are similar receive heavy edges; tracts with
dissimilar character receive light edges.
METIS's minimum-cut objective then respects community structure automatically:
cuts run along light edges (natural community transitions) and avoid heavy
edges (strong community bonds).

M.3 implements the housing character signal.
It is the third M-track implementation paper, addressing a community character
dimension that neither M.1 (economic character, workplace-side LODES data)
nor M.6 (administrative zone co-membership) capture: the residential built
environment.
How old is the housing stock?
Is it predominantly detached single-family, attached rowhouses, or large
multifamily buildings?
Do residents own or rent?
These three questions---answered by three ACS tables---define a housing
character vector that separates single-family homeowner suburbs from
apartment corridors, postwar subdivisions from pre-war rowhouse neighborhoods,
and owner-occupied historic districts from transient rental enclaves.

\paragraph{Why housing type predicts community.}
The connection between housing type and community structure is not sociological
speculation; it is encoded in law.
A single-family homeowner is, by definition, a property taxpayer with a direct
financial stake in local school quality, infrastructure maintenance, and
zoning stability.
A multifamily renter in the same tract shares none of these fiscal bonds.
The two populations have different relationships to local government, different
sensitivities to federal housing finance policy, different interests in
historic tax credit programs, and different vulnerabilities to local zoning
variance decisions~\citep{fischel2001}.
They may be geographically contiguous---a rowhouse block and an adjacent
apartment tower may literally share a property line---but they are not members
of the same community of interest in the legally relevant sense.

\paragraph{The NIMBYism signal.}
One empirically documented mechanism through which housing character becomes
a community signal is the politics of land-use opposition, commonly abbreviated
as NIMBYism (Not In My Back Yard).
Communities that are predominantly single-family and owner-occupied share
a strong revealed preference for maintaining the existing housing stock character:
they oppose upzoning, density increases, and new multifamily permitting at high
rates~\citep{fischel2001,einstein2019}.
This shared political preference is partisan-neutral---NIMBYism is documented
in liberal San Francisco and conservative Houston alike---but it is a strong
predictor of community cohesion and shared political interest.
Tracts with similar housing character (both highly single-family-and-owner)
share political interests that cross-character tracts (one single-family,
one large multifamily) do not.

\paragraph{Contribution.}
This paper makes the following contributions:
\begin{enumerate}
  \item \textbf{Four-component housing character vector} derived from three
    ACS tables (B25024, B25003, B25035), with a precise formula for each
    component and a rationale grounded in community structure theory
    (\S\ref{sec:algorithm}).
  \item \textbf{Zero-cost implementation path}: M.3 requires no new data
    download infrastructure; it extends the existing ACS demographics CSV
    with four additional columns (\S\ref{sec:algorithm}).
  \item \textbf{Empirical design}: a pre-specified experimental protocol for
    North Carolina, with falsifiable predictions about within-district housing
    homogeneity and historic neighborhood clustering (\S\ref{sec:empirical}).
  \item \textbf{Legal grounding}: a practitioner-oriented analysis of
    neighborhood stability doctrine, historic preservation district law,
    and the homeowner representation interest (\S\ref{sec:legal}).
\end{enumerate}

\paragraph{Paper organization.}
Section~\ref{sec:background} reviews the ACS methodology and the political
science literature linking housing type to community.
Section~\ref{sec:algorithm} defines the four-component character vector and
the blend formula.
Section~\ref{sec:empirical} describes the implementation evidence boundary and
future empirical design.
Section~\ref{sec:legal} provides legal grounding.
Section~\ref{sec:conclusion} concludes.
